\documentclass[conference]{IEEEtran}

\usepackage{amsmath}
\usepackage{booktabs}
\usepackage{graphicx}
\usepackage{url}

\title{An OS and Transport-Agnostic Communication Layer for Interactive Video Streaming and USB-OTG Control}

\author{
\IEEEauthorblockN{Gavril Emil-Constantin}
\IEEEauthorblockA{Technical University of "Gheorghe Asachi"\\
Iasi, Romania\\
emil-constantin.gavril@student.tuiasi.ro}
}

\begin{document}
\maketitle

\begin{abstract}
Interactive screen sharing couples latency-sensitive visual media with small,
time-critical control messages, yet most implementations bind these workloads
to one operating-system API or assume a continuously available IP network.
This paper describes the design of an operating-system- and transport-agnostic
inter-node communication layer for interactive video streaming and USB
On-The-Go (OTG) control. A shared application framing and session model is
decoupled from transport adapters for local IP, USB, and a constrained
Bluetooth fallback. Control prioritization and stale-frame expiry are design
objectives. We define an evaluation across transport performance, handover
disruption, control responsiveness, and implementation portability. The current
prototype supports continuous video streaming and transport fallback across
its implemented transports. Context-dependent automatic selection based on
inter-node payload remains future work. This draft reports no measured
performance results; the evaluation tables identify measurements for the
existing fallback and the proposed payload-aware selection.
\end{abstract}

\begin{IEEEkeywords}
heterogeneous transports, screen streaming, USB OTG, graceful degradation,
cross-platform systems, transport handover
\end{IEEEkeywords}

\section{Introduction}
Interactive screen sharing is a joint media-and-control problem. Video frames
consume most of the available bandwidth, while input and session-control
messages are small but sensitive to delay. A transport that is suitable for
one workload may be poor for the other: a high-throughput link can deliver
smooth video, whereas a low-rate nearby link may still carry useful control
traffic after the video quality has become unacceptable. Existing screen
sharing systems commonly assume a stable IP path or bind their device
integration to a particular operating system.

We investigate a communication architecture intended to preserve an
interactive session across heterogeneous nearby links. The application
protocol is separated from the link driver so that a session can use local
IP (e.g., Wi-Fi), a direct USB OTG connection, or a low-bandwidth Bluetooth
path without changing message semantics. The design treats fallback as a
media-adaptation problem as well as a connection-management problem: control
traffic is retained at high priority, while video is reduced or suspended
when the available path cannot sustain its configured rate.

This paper makes three contributions:
\begin{itemize}
  \item A transport-independent session and framing design for interactive
  video and control traffic across IP and non-IP proximity links.
  \item A candidate future-work policy for context-dependent protocol
  selection using payload class, link availability, measured performance,
  workload priorities, and hysteresis.
  \item An evaluation protocol and result templates covering latency,
  goodput, handover disruption, frame freshness, and cross-platform
  implementation cost.
\end{itemize}

These contributions distinguish implemented functionality from proposed
extension. The current prototype provides continuous video streaming and
transport fallback across the transports integrated in the system. The
payload-aware automatic protocol-selection policy described below is not yet
implemented. We therefore document the existing behavior, define an
evaluation plan, and leave quantitative results as placeholders.

\section{Problem and Design Goals}
\subsection{Communication Model}
We consider two nearby nodes, a sender and a receiver, participating in an
interactive session. The session carries three logical classes of data:
\begin{enumerate}
  \item \textit{Control}: user input, acknowledgements, capability exchange,
  and session state.
  \item \textit{Media}: encoded video frames, optionally accompanied by
  timing and key-frame metadata.
  \item \textit{Management}: discovery, transport establishment, health
  reports, and handover coordination.
\end{enumerate}
The target setting does not require an external router or cloud service.
Interfaces and transport protocols vary by endpoint; therefore, the
application layer must not depend on a particular operating-system socket,
USB, or Bluetooth API.

\subsection{Goals and Non-Goals}
The design goals are (i) consistent message semantics across transports,
(ii) bounded control latency, (iii) prevention of stale-frame queue buildup,
(iv) continuity of the session when a usable fallback exists, and (v)
platform-specific code isolated in transport adapters. This work does not
assume that all transports can carry all media at equal quality. It does not
claim bit-identical performance across operating systems, nor does the
current prototype provide production-grade authentication, encryption, or
continuous video.

\section{Architecture}
\subsection{Layer Separation}
The proposed stack has four layers:
\begin{enumerate}
  \item \textit{Application}: capture, encode/decode, and user-input handling.
  \item \textit{Session}: peer identity, capabilities, message classes,
  sequencing, timestamps, and stream lifecycle.
  \item \textit{Transport-independent framing}: bounded message headers and
  payloads, with explicit message type and delivery/expiry semantics.
  \item \textit{Transport adapters}: OS-specific implementations for local
  IP, USB OTG, and Bluetooth.
\end{enumerate}
A common adapter interface exposes establishment, send/receive, health, and
close operations. It does not expose assumptions such as an IP address to
the session layer. An adapter may internally provide a byte stream or
datagram service; framing and reassembly are selected to match that service.
This separation permits an implementation to add an adapter without
rewriting media or control logic.

Each framed message should identify its session, class, sequence number, and
generation timestamp. A video message also carries frame or stream
identification sufficient to discard obsolete data. Control messages use
bounded queues and explicit delivery semantics; media messages are subject to
deadlines. Exact field widths, authentication, and wire compatibility are
implementation decisions that must be fixed before interoperability testing.

\subsection{Transport Roles and USB OTG}
Wi-Fi/IP is the expected high-goodput path for video. USB OTG is a candidate
direct path when a device is attached as host or gadget; the adapter must
account for which endpoint controls enumeration and which USB class or
interface is available. USB availability alone does not imply a compatible
peer or an active application session. Bluetooth is treated as a constrained
fallback suitable for control and, only if measured capacity permits, a
severely reduced visual stream. The design does not assume that Bluetooth
can sustain general-purpose full-resolution video.

\begin{figure}[t]
\centering
\fbox{\parbox{0.94\columnwidth}{\centering
\textbf{Architecture figure placeholder}\\[2pt]
Application capture / input
$\rightarrow$ session framing and scheduler
$\rightarrow$ IP, USB OTG, or Bluetooth adapters\\
Show peer capability exchange, link-health feedback, and fallback control
path in the final figure.}}
\caption{Proposed separation of application traffic from heterogeneous link adapters.}
\label{fig:architecture}
\end{figure}

\section{Scheduling and Fallback}
\subsection{Traffic Treatment}
Control messages are scheduled ahead of bulk media, with queue limits to
avoid head-of-line delay caused by large frame transfers. For each video
frame, the sender attaches a capture timestamp and discards frames whose
deadline has passed before transmission. Where the encoder supports it,
adaptation may reduce bitrate, resolution, or frame rate; if no useful
visual rate can be sustained, the system continues control-only operation
and exposes the degraded state to the user.

\subsection{Candidate Selection Policy}
Let $T$ be the set of available transports. For transport $i$, estimate
round-trip time $L_i$, jitter $J_i$, effective goodput $B_i$, loss or
retransmission rate $P_i$, energy cost $E_i$, and transition cost $S_i$.
A candidate score for workload $w$ is
\begin{equation}
C_i(w) = \alpha_w L_i + \beta_w J_i +
\frac{\gamma_w}{B_i+\epsilon} + \delta_w P_i +
\eta_w E_i + \kappa_w S_i ,
\label{eq:cost}
\end{equation}
where weights depend on whether the session prioritizes interactive control,
visual quality, or energy use. A candidate is eligible only if it meets
minimum workload-specific constraints, such as a minimum control service
rate. The score is a design hypothesis; its weights and thresholds require
empirical calibration and are not yet selected.

The existing fallback mechanism handles loss or degradation of the active
transport. The context-dependent selector described here is future work: it
would select a protocol using the inter-node payload class and measured link
conditions rather than connectivity alone. In that design, it would not
switch on a single transient sample; it would require sustained degradation
before switching down and a sustained improvement margin before switching
back. A handover would retain session state and control sequence numbers while
discarding late data from the old path. The current implementation's exact
fallback triggers, recovery behavior, and supported transport pairs should
be recorded with the experiment configuration.

\section{Prototype Status}
The current system combines a Flutter interface with a Rust communication
backend and supports continuous video streaming and transport fallback.
Fallback and streaming are treated as implemented capabilities; their
performance and coverage across operating systems and transport pairs remain
to be quantified in the evaluation. The system does not yet perform
context-dependent automatic protocol selection based on inter-node payload
type (for example, distinguishing control messages from video when choosing
a transport). That selector, including its policy thresholds and adaptation
behavior, is proposed future work rather than a measured feature of the
current prototype.

\section{Experimental Methodology}
\subsection{Questions}
The evaluation should answer:
\begin{itemize}
  \item \textbf{RQ1:} How do latency, jitter, goodput, and frame freshness
  differ across the supported links and tested operating systems?
  \item \textbf{RQ2:} What interruption and loss of control responsiveness
  occur during fallback and recovery?
  \item \textbf{RQ3:} Does traffic prioritization preserve control
  responsiveness when video saturates a constrained link?
  \item \textbf{RQ4:} Under an explicit energy budget, can the proposed
  energy-aware selector reduce device energy use while preserving control
  service, compared with the existing fallback?
  \item \textbf{RQ5:} What portion of the implementation is shared across
  operating systems, and what runtime/resource costs accompany the design?
\end{itemize}

\subsection{Setup and Procedure}
Use at least one Android device and one desktop endpoint for each supported
platform, documenting model, OS version, radio configuration, USB role,
codec, resolution, and power mode. Report direct-link and baseline conditions
separately. For each transport, run a control-only workload, a fixed-rate
video workload, and a mixed interactive workload. Repeat trials across
multiple link conditions and report sample counts, median, 95th percentile,
dispersion, and confidence intervals. Randomize run order where practical.

For fallback-delay tests, trigger the conditions handled by the existing
mechanism, such as abrupt loss of the active link. Separately test gradual
degradation below a configured service threshold for the proposed selector.
Timestamp (i) the trigger, (ii) degradation detection, (iii) selection of the replacement
transport, (iv) readiness of that transport, (v) delivery of the first
control message, and (vi) display of the first valid, non-stale video frame.
Report each interval separately as well as total control interruption and
video interruption; this distinguishes slow detection from slow reconnection
or media recovery. Test each supported ordered transport pair, including
fallback and recovery directions, and record failed transitions, lost or
expired frames, and control-message delay. Compare the existing fallback with a fixed-link baseline. Evaluate the
future payload-aware selector separately, against the existing fallback and,
if implemented, a no-hysteresis policy.

For evaluation of the proposed energy-aware selector, define a device energy reserve or remaining
runtime budget before each trial, then exercise the selector at several
reserve levels under matched control-only and mixed video/control workloads.
Record the state and threshold that caused a decision, selected transport,
quality reduction, and any time spent in control-only mode. Measure device
power with an external power meter where possible; otherwise use a documented
platform energy estimate and identify it as an estimate. Subtract or report
idle power separately, and normalize energy by session duration and delivered
useful frames. Compare the energy-aware policy against the same session with
energy weighting disabled. A lower-power transport must be established by
measurement rather than assumed from its link type; report control latency
and frame freshness alongside energy savings to expose quality trade-offs.
Keep device, battery state, workload, and radio conditions matched across
paired runs.

\subsection{Metrics}
Measure end-to-end control RTT and one-way delay where clock synchronization
allows; video goodput; frame age at display; delivered and expired frame
counts; frame rate; fallback detection, selection, readiness, and recovery
delays; control and video interruption; reconnection success; CPU, memory,
and energy; and shared versus platform-specific source lines. For energy
trials report measurement equipment or estimation method, idle baseline,
incremental energy, and energy per useful delivered frame. State the
capture-to-display measurement method and clock error. Network RTT must not
be presented as capture-to-display latency.

\section{Results}
\textit{No experiment results are available in this draft. Replace every
placeholder with measured values and the specified test conditions before
submission.}

\begin{table}[t]
\caption{Per-Transport Performance (Placeholder)}
\label{tab:transport-results}
\centering
\begin{tabular}{lccccc}
\toprule
Transport & Goodput & Control RTT & Jitter & Frame age & Energy\\
 & (Mb/s) & (ms) & (ms) & (ms) & (mWh/min)\\
\midrule
Wi-Fi & TBD & TBD & TBD & TBD & TBD\\
USB OTG & TBD & TBD & TBD & TBD & TBD\\
Bluetooth & TBD & TBD & TBD & TBD & TBD\\
\bottomrule
\end{tabular}
\end{table}

\begin{table}[t]
\caption{Fallback and Recovery Delay (Placeholder)}
\label{tab:handover-results}
\centering
\scriptsize
\begin{tabular}{lccccc}
\toprule
Transition & Detect & Select & Ready & Control & Video\\
 & (ms) & (ms) & (ms) & gap (ms) & gap (ms)\\
\midrule
Wi-Fi $\rightarrow$ USB & TBD & TBD & TBD & TBD & TBD\\
USB $\rightarrow$ Bluetooth & TBD & TBD & TBD & TBD & TBD\\
Wi-Fi $\rightarrow$ Bluetooth & TBD & TBD & TBD & TBD & TBD\\
\bottomrule
\end{tabular}
\end{table}

\begin{table}[t]
\caption{Energy-Constrained Fallback (Placeholder)}
\label{tab:energy-fallback-results}
\centering
\scriptsize
\begin{tabular}{lccccc}
\toprule
Policy & Reserve & Selected & Energy & Control & Frame\\
 & (Wh or \%) & link & saving (\%) & RTT (ms) & age (ms)\\
\midrule
Energy-aware & TBD & TBD & TBD & TBD & TBD\\
Energy weight off & TBD & TBD & TBD & TBD & TBD\\
\bottomrule
\end{tabular}
\end{table}

\begin{figure}[t]
\centering
\fbox{\parbox{0.94\columnwidth}{\centering
\textbf{Results plot placeholder}\\[2pt]
Plot control RTT and displayed frame age over time for a forced
Wi-Fi-to-fallback transition. Mark degradation trigger, detection,
replacement-link readiness, first control delivery, and first valid frame.}}
\caption{Suggested timeline for measuring fallback delay and interruption.}
\label{fig:handover}
\end{figure}

\section{Discussion}
The main design trade-off is continuity versus quality: a fallback may keep
control available while providing little or no useful video. Therefore,
success should not be defined only as socket/session survival. Evaluation
must separately report control continuity, frame freshness, and visual
quality. In particular, reliable retransmission of expired video can worsen
interactive performance; policy should be selected by message class rather
than by a single transport-wide reliability setting.

Transport portability also does not imply equivalent platform capabilities.
USB OTG role negotiation and device permissions differ across operating
systems, while Bluetooth profiles and throughput vary by device. These
differences must be reported as limitations, not hidden behind the common
adapter interface. Security is another prerequisite: the current Wi-Fi
prototype is intended for trusted LAN testing and lacks production
authentication and encryption. A deployable system must authenticate peers
and protect control and media traffic independently of link type.

\section{Related Work}
Screen-sharing studies compare native and clientless approaches and motivate
measuring end-to-end behavior rather than inferring user experience from
transport throughput alone \cite{huang2021}. Handover research has examined
predicting Wi-Fi loss and preparing alternate paths using multipath transport
\cite{hochst2019}. Such work informs link monitoring and continuity, but the
present design focuses on application-level interactive media and control
across IP and non-IP proximity transports. Media-independent handover
standards provide a broader framework for separating link events from
higher-layer decisions \cite{ieee80221}.

The contribution sought here is not a new codec or a claim that Rust or
cross-platform UI frameworks are novel. It is the integration and evaluation
of a transport-independent interactive session model, traffic-aware
degradation, and fallback across distinct local link types. This claim
remains to be tested against implemented systems and measured baselines.

\section{Limitations and Future Work}
The current system supports continuous video streaming and transport
fallback, but this draft contains no measured performance, energy,
reliability, or code-sharing results. The principal future-work item is
context-dependent automatic protocol selection using inter-node payload
characteristics together with measured link conditions. Future evaluation
must compare this selector with the existing fallback, quantify its effects
on control delay, frame freshness, and energy use, and document the supported
transport/operating-system combinations. Broader multi-node routing and
battery-aware topology selection are outside the scope of this pairwise
study.

\section{Conclusion}
We presented a design for inter-node interactive streaming that separates
session semantics from operating-system-specific transport adapters and
gives control and video distinct scheduling and degradation behavior.
The current prototype provides continuous video streaming and transport
fallback. Context-dependent automatic protocol selection based on inter-node
payload remains future work, and the performance and energy benefits of both
the existing fallback and that proposed extension must be established
empirically.

\begin{thebibliography}{00}

\bibitem{huang2021}
C.-Y. Huang, Y.-C. Cheng, G.-Z. Huang, C.-L. Fan, and C.-H. Hsu,
``On the performance comparisons of native and clientless real-time screen
sharing technologies,'' \emph{ACM Transactions on Multimedia Computing,
Communications, and Applications}, vol. 17, no. 2, Art. no. 54, 2021,
doi: 10.1145/3437881.

\bibitem{hochst2019}
J. H{\"o}chst \emph{et al.}, ``Learning Wi-Fi connection loss predictions
for seamless vertical handovers using Multipath TCP,'' in
\emph{Proc. IEEE 44th Conference on Local Computer Networks (LCN)}, 2019,
doi: 10.1109/LCN44214.2019.8990753.

\bibitem{ieee80221}
IEEE, \emph{IEEE Standard for Local and Metropolitan Area Networks---Media
Independent Handover Services}, IEEE Std 802.21-2008, 2009.

\end{thebibliography}

\end{document}
